\section{Discussion}
The idea of test-time training also makes sense for other tasks, such as segmentation and detection, and in other fields, such as speech recognition and natural language processing. 
For machine learning practitioners with prior domain knowledge in their respective fields, their expertise can be leveraged to design better special-purpose self-supervised tasks for test-time training.
Researchers for general-purpose self-supervised tasks can also use  test-time training as an evaluation benchmark, in addition to the currently prevalent benchmark of pre-training and fine-tuning.

More generally, we hope this paper can encourage researchers to abandon the self-imposed constraint of a fixed decision boundary for testing, or even the artificial division between training and testing altogether. 
Our work is but a small step toward a new paradigm where much of the learning happens {\em after} a model is deployed.
